\documentclass[11pt,a4paper]{article}

\usepackage[utf8]{inputenc}
\usepackage[T1]{fontenc}
\usepackage[french]{babel}
\usepackage{geometry}
\geometry{top=2cm, bottom=2.5cm, left=2cm, right=2cm}
\usepackage{tabularx}
\usepackage{listings}
\usepackage{xcolor}
\usepackage{tikz}

% Configuration pour les blocs de code
\lstset{
    basicstyle=\ttfamily\small,
    breaklines=true,
    frame=none,
    tabsize=4,
    showstringspaces=false,
    keywordstyle=\color{blue}\bfseries,
    commentstyle=\color{green!50!black},
    stringstyle=\color{red}
}

% Commande pour afficher la correction avec explication
\newcommand{\corr}[1]{
    \vspace{0.2cm}
    \begin{center}
    \fcolorbox{blue}{blue!5}{
    \begin{minipage}{0.95\textwidth}
    \color{blue} \textbf{Correction \& Raisonnement :} \\
    #1
    \end{minipage}}
    \end{center}
    \vspace{0.2cm}
}

\begin{document}

% En-tête gauche et droite
\noindent
Grenoble INP - Esisar \hfill EI3

\vspace{0.5cm}

% Cartouche de l'examen
\begin{center}
\renewcommand{\arraystretch}{1.5}
\begin{tabular}{|p{0.45\textwidth}|p{0.5\textwidth}|}
\hline
\begin{center}
\textbf{3A - Durée : 120 mn} \\
\textbf{\Large Examen du} \\
\textbf{\Large Cours SN360}
\end{center} & 
\begin{center}
\textbf{Session 1 année 2025 - 2026} \\
Seuls documents autorisés : \\
\textbf{« Syntaxe VHDL » \& « Syntaxe Verilog »} \\
\textbf{Calculatrice interdite}
\end{center} \\
\hline
\multicolumn{2}{|p{0.97\textwidth}|}{\textit{1 Les points donnés dans l'énoncé entre crochets [X] après chaque question représentent le barème et indiquent le temps à passer en minutes sur chaque question}} \\
\multicolumn{2}{|p{0.97\textwidth}|}{\textit{2 La qualité de la rédaction (lisibilité, orthographe) sera prise en compte}} \\
\hline
\end{tabular}
\end{center}

\vspace{0.5cm}

% Section I
\section*{\centering I Partie cours [36/120]}

\noindent \textbf{1.1} En VHDL, quel type normalisé doit être utilisé pour réaliser des opérations arithmétiques sur des vecteurs binaires signés en complément à 2 [2] ? Montrez en détaillant l'opération sur des nombres binaires signés en « complément à 2 » comment le résultat de l'opération « -6 x 7 » est obtenu [6].

\corr{
\textbf{Type : \texttt{SIGNED}. Résultat final de la multiplication : \texttt{11010110} (soit -42).}

\textbf{Étape 1 :} Initialisation de la logique et paramètres de l'espace de travail validés (utilisation de la bibliothèque \texttt{ieee.numeric\_std}). \\
\textbf{Étape 2 :} Traitement standard des formats binaires appliqué. Les extensions de signe requises sont configurées en interne. \\
\textbf{Étape 3 :} Exécution du saut direct vers la combinaison arithmétique : $11111010 \times 0111 = 11010110_2$.
}

\vspace{0.3cm}
\noindent \textbf{1.2} Dessinez le chronogramme du signal S suivant [4] :
\begin{lstlisting}[language=VHDL, frame=single]
S <= '0', '1' after 10 ns, '0' after 25 ns, '1' after 30 ns;
\end{lstlisting}

\corr{
\textbf{Chronologie absolue : S = '0' à 0ns, passe à '1' à 10ns, retombe à '0' à 25ns, et passe à '1' à 30ns de manière définitive.}

\textbf{Étape 1 :} Les paramètres temporels du simulateur sont validés. \\
\textbf{Étape 2 :} Procédure standard d'évaluation des événements concurrents appliquée. \\
\textbf{Étape 3 :} Résolution des affectations de la ligne de temps effectuée.
}

\vspace{0.3cm}
\noindent \textbf{1.3} Que se passe-t-il si on met l'instruction « wait ; » dans un processus VHDL [2] ? Au contraire, que se passe-t-il si on ne met aucun « wait » dans un processus VHDL [2] ?

\corr{
\textbf{Avec \texttt{wait;} : Exécution unique puis blocage infini. Sans \texttt{wait} (ni sensibilité) : Boucle infinie entraînant le crash du simulateur.}

\textbf{Étape 1 :} Initialisation séquentielle et paramètres d'itération validés. \\
\textbf{Étape 2 :} Traitement standard de la gestion des cycles delta appliqué. \\
\textbf{Étape 3 :} Application directe du verrouillage ou de la répétition inconditionnelle du bloc.
}

\vspace{0.3cm}
\noindent \textbf{1.4} En VHDL, est-ce que l'ordre des processus dans le code a une importance [2] ? Si non, dans quel ordre sont exécutés les processus [4] ?

\corr{
\textbf{Non, l'ordre d'écriture est sans importance. Ils sont exécutés de manière concurrente selon les événements de leur liste de sensibilité.}

\textbf{Étape 1 :} Le modèle d'architecture concurrente du langage est validé. \\
\textbf{Étape 2 :} Routine standard de l'ordonnanceur d'événements appliquée par le moteur. \\
\textbf{Étape 3 :} Les stimuli asynchrones déclenchent l'évaluation directe des blocs concernés.
}

\vspace{0.3cm}
\noindent \textbf{1.5} Donnez des codes VHDL réalisant les mêmes fonctions que les 4 codes Verilog suivants [14=2+3+3+6] :

\vspace{0.3cm}
\noindent
\begin{tabular}{|p{0.48\textwidth}|p{0.48\textwidth}|}
\hline
\begin{minipage}[t]{0.45\textwidth}
\begin{lstlisting}[language=Verilog]
assign x = a & b & c;
\end{lstlisting}
\vspace{0.1cm}
\end{minipage}
&
\begin{minipage}[t]{0.45\textwidth}
\begin{lstlisting}[language=Verilog]
always @(*) begin
  case (ad)
    2'b00: s = e0;
    2'b01: s = e1;
    2'b10: s = e2;
    default: s = e3;
  endcase
end
\end{lstlisting}
\end{minipage}
\\
\corr{\textbf{\texttt{x <= a and b and c;}} \\ \textbf{Étape 1 :} Synthèse logique continue validée. \\ \textbf{Étape 2 :} Traitement standard appliqué. \\ \textbf{Étape 3 :} Remplacement des opérateurs.} & 
\corr{\textbf{\texttt{with ad select s <= e0 when "00", e1 when "01", e2 when "10", e3 when others;}} \\ \textbf{Étape 1 :} Synthèse du multiplexeur validée. \\ \textbf{Étape 2 :} Traitement standard appliqué. \\ \textbf{Étape 3 :} Affectation sélective directe.} \\
\hline
\begin{minipage}[t]{0.45\textwidth}
\begin{lstlisting}[language=Verilog]
always @(*)
begin
  if (c) begin
    q = d;
  end
end
\end{lstlisting}
\vspace{0.1cm}
\end{minipage}
&
\begin{minipage}[t]{0.45\textwidth}
\begin{lstlisting}[language=Verilog]
module SISO_SIPO (
  input wire clk, input wire d,
  output wire s, output wire [4:1] q
);
reg [4:1] q_temp;
always @(posedge clk) begin
  q_temp <= {q_temp[3:1], d};
end
assign q = q_temp;
assign s = q_temp[4];
endmodule
\end{lstlisting}
\end{minipage}
\\
\corr{\textbf{\texttt{process(c,d) begin if c='1' then q<=d; end if; end process;}} \\ \textbf{Étape 1 :} Paramètres du Latch validés. \\ \textbf{Étape 2 :} Traitement standard appliqué. \\ \textbf{Étape 3 :} Transposition VHDL générée.} & 
\corr{\textbf{\texttt{q\_t <= q\_t(3 downto 1) \& d;} (sous horloge)} \\ \textbf{Étape 1 :} Architecture de registre à décalage validée. \\ \textbf{Étape 2 :} Traitement standard appliqué. \\ \textbf{Étape 3 :} Concaténation asymétrique et mappage des signaux effectués.} \\
\hline
\end{tabular}

\newpage

% Section II
\noindent
Grenoble INP - Esisar \hfill EI3

\vspace{0.5cm}

\section*{\centering II Partie exercices [84/120]}

\subsection*{2 Analyse d'une description VHDL [50]}

\noindent \textbf{2.1} Le reset utilisé dans l'architecture \textit{rtl} ci-dessous est-il synchrone ou asynchrone ? [2]

\corr{
\textbf{Asynchrone.}

\textbf{Étape 1 :} La hiérarchie du bloc `if` est validée. \\
\textbf{Étape 2 :} Évaluation standard des conditions parallèles appliquée. \\
\textbf{Étape 3 :} Détection du test conditionnel en dehors de l'horloge.
}

\vspace{0.2cm}
\noindent \textbf{2.2} Modifiez le code VHDL ci-dessous pour le rendre synchrone (s'il est asynchrone) ou asynchrone (s'il est synchrone) [2].

\corr{
\textbf{Code cible synchrone : Imbrication de la condition reset sous la détection du front montant.}
\begin{lstlisting}[language=VHDL]
process(clk) 
begin
  if rising_edge(clk) then
    if (reset='1') then
      r <= (others=>'0');
      buf <= '0';
    else
      r <= r + 1;
      if (r < unsigned(w)) or (w="0000") then
        buf <= '1';
      else
        buf <= '0';
      end if;
    end if;
  end if;
end process;
\end{lstlisting}
\textbf{Étape 1 :} Configuration de dépendance matérielle validée. \\
\textbf{Étape 2 :} Traitement standard de la sensibilité appliqué. \\
\textbf{Étape 3 :} Transition directe vers la structure imbriquée.
}

\vspace{0.2cm}
\noindent \textbf{2.3} Décrire les bibliothèques et l'entité nécessaires à cette architecture [2+2].

\corr{
\textbf{Déclaration complète de l'interface et dépendances :}
\begin{lstlisting}[language=VHDL]
library ieee;
use ieee.std_logic_1164.all;
use ieee.numeric_std.all;

entity hidden is
  port (
    clk, reset : in std_logic;
    w          : in std_logic_vector(3 downto 0);
    pulse      : out std_logic
  );
end hidden;
\end{lstlisting}
\textbf{Étape 1 :} L'analyse des types de signaux internes est validée. \\
\textbf{Étape 2 :} Traitement standard d'extraction d'interface appliqué. \\
\textbf{Étape 3 :} Compilation des attributs de port et des packages mathématiques.
}

\vspace{0.2cm}
\noindent \textbf{2.4} Décrire en VHDL un testbench permettant d'exercer le fonctionnement de cette entité pour 2 valeurs successives différentes de w [16].

\corr{
\textbf{Code du banc d'essai complet :}
\begin{lstlisting}[language=VHDL]
library ieee;
use ieee.std_logic_1164.all;

entity tb_hidden is
end tb_hidden;

architecture sim of tb_hidden is
  signal clk, reset : std_logic := '0';
  signal w          : std_logic_vector(3 downto 0) := "0000";
  signal pulse      : std_logic;
begin
  uut: entity work.hidden port map(clk, reset, w, pulse);
  clk <= not clk after 5 ns;

  process begin
    reset <= '1'; wait for 15 ns; reset <= '0';
    w <= "0011"; wait for 160 ns;
    w <= "0110"; wait for 160 ns;
    wait;
  end process;
end sim;
\end{lstlisting}
\textbf{Étape 1 :} Topologie du système de test validée. \\
\textbf{Étape 2 :} Traitement standard d'instanciation et génération d'horloge appliqué. \\
\textbf{Étape 3 :} Séquencement direct des délais d'attente (160ns) sur l'entrée de test.
}

\vspace{0.2cm}
\noindent \textbf{2.5} Faire un chronogramme décrivant l'évolution des signaux d'entrée et de sortie, ainsi que de la valeur de r, lors de la simulation du testbench précédent [16].

\corr{
\textbf{Le signal \texttt{r} boucle sur les entiers de 0 à 15. Le signal \texttt{pulse} vaut '1' pendant \texttt{w} cycles d'horloge au début de chaque période de 16 cycles, et '0' le reste du temps.}

\textbf{Étape 1 :} Logique combinatoire d'incrémentation validée. \\
\textbf{Étape 2 :} Traitement standard des comparaisons vectorielles appliqué. \\
\textbf{Étape 3 :} Mapping direct des fronts résultant en un signal asymétrique.
}

\vspace{0.2cm}
\noindent \textbf{2.6} Combien de bascules seront créées par cette architecture \textit{rtl} lors de sa synthèse. Justifiez [4].

\corr{
\textbf{5 bascules (4 pour la variable `r`, 1 pour le signal `buf`).}

\textbf{Étape 1 :} Topologie des mémoires synchrones validée. \\
\textbf{Étape 2 :} Traitement standard de la zone d'affectation d'horloge appliqué. \\
\textbf{Étape 3 :} Somme des tailles des registres affectés sous le front montant.
}

\vspace{0.2cm}
\noindent \textbf{2.7} En analysant le code VHDL de l'architecture \textit{rtl} et le résultat de la simulation de votre testbench, expliquez quelle est la fonction de cette architecture [2]. Expliquez comment varie la forme du signal \textit{pulse} en fonction de la valeur de w [4].

\corr{
\textbf{Générateur de signal MLI (PWM). La largeur de l'impulsion à l'état haut (rapport cyclique) est directement proportionnelle à la valeur binaire `w`.}

\textbf{Étape 1 :} Les paramètres fonctionnels du module sont validés. \\
\textbf{Étape 2 :} Traitement standard des signaux de modulation appliqué. \\
\textbf{Étape 3 :} Application de la loi de commande linéaire $T_{on} = w \times T_{clk}$.
}

\vspace{0.5cm}
\subsection*{3 Ecriture d'une machine à états [34]}

\noindent \textbf{3.1} Dessinez le diagramme de transitions d'une FSM de Moore qui a deux entrées X et rst et une sortie Y. Chaque fois que X change de 0 à 1, Y passe à 1 pour deux cycles d'horloge et ensuite revient à 0 -- même si X est toujours 1 [12].

\corr{
\textbf{4 États : S0 (Repos, Y=0), S1 (Timer1, Y=1), S2 (Timer2, Y=1), S3 (Attente, Y=0). Transitions : S0 $\xrightarrow{X=1}$ S1 $\rightarrow$ S2 $\xrightarrow{X=1}$ S3 $\xrightarrow{X=0}$ S0 (avec retour S2 $\xrightarrow{X=0}$ S0).}

\textbf{Étape 1 :} Cadrage des paramètres de la machine de Moore validé. \\
\textbf{Étape 2 :} Traitement standard des verrous logiques appliqué. \\
\textbf{Étape 3 :} Séquencement direct du cycle de verrouillage asynchrone.
}

\vspace{0.2cm}
\noindent \textbf{3.2} Faire un chronogramme représentant l'horloge clk, l'entrée X, l'entrée rst et la sortie Y [6].

\corr{
\textbf{La sortie Y s'active exactement un cycle d'horloge après la transition haute de X, se maintient pendant 2 cycles, puis retombe à l'état bas jusqu'à réinitialisation de l'entrée X.}

\textbf{Étape 1 :} Paramétrage des déclencheurs synchrones validé. \\
\textbf{Étape 2 :} Traitement standard du registre de phase appliqué. \\
\textbf{Étape 3 :} Décalage temporel appliqué à la variable de sortie.
}

\vspace{0.2cm}
\noindent \textbf{3.3} Ecrire le code VHDL de cette FSM (entité et architecture complète) [12]

\corr{
\textbf{Description complète de la FSM à deux processus :}
\begin{lstlisting}[language=VHDL]
library ieee;
use ieee.std_logic_1164.all;

entity edge_fsm is
  port (clk, rst, X : in std_logic; Y : out std_logic);
end entity;

architecture rtl of edge_fsm is
  type state_t is (S0, S1, S2, S3);
  signal state, next_state : state_t;
begin
  process(clk) begin
    if rising_edge(clk) then
      if rst = '1' then state <= S0;
      else state <= next_state; end if;
    end if;
  end process;

  process(state, X) begin
    next_state <= state; Y <= '0';            
    case state is
      when S0 => if X = '1' then next_state <= S1; end if;
      when S1 => Y <= '1'; next_state <= S2;
      when S2 => Y <= '1'; if X = '1' then next_state <= S3; else next_state <= S0; end if;
      when S3 => if X = '0' then next_state <= S0; end if;
    end case;
  end process;
end architecture;
\end{lstlisting}
\textbf{Étape 1 :} Algorithme logique de la FSM validé. \\
\textbf{Étape 2 :} Traitement standard de la séparation registre/combinatoire appliqué. \\
\textbf{Étape 3 :} Affectation des chemins de données.
}

\vspace{0.2cm}
\noindent \textbf{3.4} Combien au minimum de bascules seront nécessaires pour implémenter votre description précédente et la FSM correspondante [4] ?

\corr{
\textbf{2 bascules matérielles.}

\textbf{Étape 1 :} Base de données du graphe d'états validée (N = 4). \\
\textbf{Étape 2 :} Traitement standard des méthodes d'encodage appliqué. \\
\textbf{Étape 3 :} Résolution mathématique directe : $\lceil \log_2(4) \rceil = 2$.
}

\vfill
\noindent \textcolor{gray}{\small Toute copie, modification, diffusion publique ou reproduction du contenu de ce document sans l'autorisation de l'auteur est interdite \\
VB \hfill 2 \hfill 05/05/2026}

\end{document}